% answers_short.tex —— 1977 年·江苏徐州市、徐州地区初试卷 · 简短答案速查
% 由 tools/make_short_answers.py 自 answers.tex 的 \gksolhead/\gkanswer 提取，
% 并与 1977/1977_数学真题_江苏徐州市,徐州地区初试_答案要点.md 逐题比对后生成，勿手改（改源 MD 后重跑）。

\begin{gknotice}
\textbf{使用说明}：本节只给最终答案便于快速核对，完整推导见解析版。答案为编写者独立推导，\textbf{非官方标准答案}。
\end{gknotice}

\gkpart{题目}
\begin{itemize}[leftmargin=2.6em]
  \item[\textbf{1.}] $6$。
  \item[\textbf{2.}] (1) $-(\sqrt{3}+2)$；(2) $2$
  \item[\textbf{3.}] (1) $-\dfrac{\sqrt{3}}{2}$；(2) $-\dfrac{\sqrt{3}}{3}$
  \item[\textbf{4.}] $x=2$，$y=-3$。
  \item[\textbf{5.}] $\dfrac{2}{x^2+2x-3}$。
  \item[\textbf{6.}] $BC=9\sqrt{2}$ 海浬（$\approx12.7$ 海浬）。
  \item[\textbf{7.}] 宽取 $5$ 米（此时长为 $10$ 米），最大面积 $50$ 平方米。
  \item[\textbf{8.}] $\lg1=0$，$\lg2=0.3010$，$\lg3=0.4771$，$\lg4=0.6020$，$\lg5=0.6990$，$\lg6=0.7781$，$\lg8=0.9030$，$\lg9=0.9542$，$\lg10=1$。
  \item[\textbf{9.}] (1) $x^2=-18y$；(2) 见解析版；(3) 焦点 $\left(0,-\dfrac{9}{2}\right)$，准线 $y=\dfrac{9}{2}$
  \item[\textbf{10.}] $a=3$，$b=27$；或 $a=27$，$b=3$。
\end{itemize}
